\chapter{Proteger o sistema nervoso}\label{ch:g8:protecting-the-nervous-system}

Dois capítulos mostraram o que o \omterm{def:g5:brain-in-command:system}{sistema nervoso} é: uma fiação cujas
junções aprendem, decidem --- e podem ser enganadas
(\cref{rem:g8:nervous-communication:vulnerable}). O encerramento deste
ano transforma essa anatomia em proteção. O que a falta de \omterm{prop:g1:healthy-habits:needs}{sono}, o
barulho, as telas e as substâncias psicoativas de fato fazem, junção
por junção, com a máquina que é você --- e como é, na prática, proteger
essa máquina aos catorze anos?

\section{Manutenção: o sono e os limites dos sentidos}

\begin{proposition}[O sono é manutenção de sinapses]\label{prop:g8:protecting-the-nervous-system:sleep}
O arquivamento e os reparos da
\cref{prop:g5:brain-in-command:protect} têm uma leitura no nível das
junções: no \omterm{prop:g1:healthy-habits:needs}{sono}, as \omterm{def:g8:nervous-communication:synapse}{sinapses} fortalecidas do dia são consolidadas ---
a prática que se torna permanente
(\cref{ex:g8:nervous-communication:juggler}) acontece em boa parte
durante a noite --- e a faxina do \omterm{def:g5:brain-in-command:system}{cérebro} limpa o desgaste químico do
dia. Noites curtas deixam o aprendizado de ontem arquivado pela metade
e a pesagem de amanhã lenta: o \omterm{def:g5:brain-in-command:system}{cérebro} \omterm{def:g3:stages-of-life:stages}{adolescente}, em plena religação
(\cref{rem:g8:puberty:moods}) e com o relógio escorregado para tarde,
precisa das horas dele justamente quando a escola começa mais cedo.
\end{proposition}
\admitted

\begin{example}[O paradoxo da revisão]\label{ex:g8:protecting-the-nervous-system:revision}
Dois estudantes antes de uma prova: um revisa até a meia-noite e segue
até as duas; o outro revisa até as dez e dorme. O segundo em geral
ganha --- a matéria atravessou as junções dele \emph{e foi
consolidada}; as horas extras do primeiro custaram justamente o
arquivamento que faz a revisão pegar. ``Dormir sobre o assunto'' é
engenharia de \omterm{def:g8:nervous-communication:synapse}{sinapse}.
\end{example}

\begin{proposition}[Volume alto é dose]\label{prop:g8:protecting-the-nervous-system:noise}
As \omterm{def:g5:first-look-at-cells:cell}{células} sensoriais do \omterm{def:g1:the-five-senses:sense}{ouvido} --- os receptores finos como fio de
cabelo que transformam som em \omterm{def:g8:nervous-communication:neuron}{mensagens nervosas} --- não voltam a
crescer (a \cref{prop:g1:the-five-senses:fragile} disse isso com
delicadeza; aqui está o mecanismo). O som alto mata essas \omterm{def:g5:first-look-at-cells:cell}{células} por
dose: volume vezes exposição. Shows e fones no volume máximo gastam,
hora após hora, um capital de receptores que \omterm{prop:g1:healthy-habits:needs}{sono} nenhum devolve --- o
apito depois de uma noite barulhenta
(\cref{ex:g1:the-five-senses:loud}) é o aviso de dose excessiva. A
regra é um orçamento: mais baixo, mais curto, com pausas --- e protetor
auricular onde os amplificadores mandam.
\end{proposition}
\admitted

\section{As substâncias que enganam as junções}

\begin{proposition}[As substâncias psicoativas agem nas sinapses]\label{prop:g8:protecting-the-nervous-system:substances}
Toda substância que muda o humor, a percepção ou o estado de alerta
funciona nas fendas da
\cref{def:g8:nervous-communication:synapse} --- imitando um
mensageiro, bloqueando um ou forçando a liberação dele:
\begin{itemize}
  \item o \emph{álcool} amortece as junções do sistema inteiro: o
        circuito fica lento de ponta a ponta --- julgamento, \omterm{def:g5:brain-in-command:reaction}{tempo de
        reação} (os metros do \cref{ex:g5:brain-in-command:driver}
        aumentam), equilíbrio, memória; as junções ainda em construção
        do \omterm{def:g5:brain-in-command:system}{cérebro} \omterm{def:g3:stages-of-life:stages}{adolescente} recebem esse amortecimento mais fundo;
  \item a \emph{nicotina do tabaco} imita um mensageiro natural e,
        repetida, religa as junções de recompensa rumo à
        \emph{\omterm{prop:g5:healthy-choices:substances}{dependência}} (a armadilha da
        \cref{prop:g5:healthy-choices:substances}, agora com o
        mecanismo dela: as \omterm{def:g8:nervous-communication:synapse}{sinapses} se ajustam a esperar a dose --- e
        protestam sem ela);
  \item a \emph{maconha} imita outra família de mensageiros,
        embaralhando justamente as junções da memória, da motivação e
        do senso de tempo --- de novo, mais cara num \omterm{def:g5:brain-in-command:system}{cérebro} em obras;
  \item quanto mais pesada a substância, mais grosseira a forçada nas
        junções --- e mais rápida a religação de recompensa que é o
        motor da \omterm{prop:g5:healthy-choices:substances}{dependência}.
\end{itemize}
\end{proposition}
\admitted

\begin{remark}[A dependência, mecanicamente]\label{rem:g8:protecting-the-nervous-system:addiction}
A armadilha da \cref{rem:g5:healthy-choices:never} agora tem um
diagrama de fiação: as junções de recompensa, dosadas repetidas vezes,
\emph{se adaptam} --- receptores reduzidos, limiares elevados --- até a
substância ser necessária para a pessoa se sentir normal, e a ausência
dela virar um protesto do sistema. Essa \omterm{def:g4:adapted-to-their-world:adaptation}{adaptação} é a assimetria de
entrada fácil e saída difícil feita de \omterm{def:g8:nervous-communication:synapse}{sinapses}: começar é uma
decisão; parar é lutar contra junções religadas. O \omterm{def:g5:brain-in-command:system}{cérebro}
\omterm{def:g3:stages-of-life:stages}{adolescente}, ainda ajustável, se religa mais rápido --- e é por isso
que toda estatística encontra quem começou cedo preso mais fundo, e é
por isso que nunca começar continua sendo o movimento barato.
\end{remark}

\begin{example}[O ciclo da tela]\label{ex:g8:protecting-the-nervous-system:screens}
Sem molécula nenhuma, e nas mesmas junções: os feeds e os jogos são
ajustados --- de propósito --- para cutucar a fiação de recompensa com
prêmios pequenos e imprevisíveis, o esquema que treina as junções com
mais força. O resultado é um primo brando da armadilha: uma checagem
que resiste a parar, noites que somem, \omterm{prop:g1:healthy-habits:needs}{sono} empurrado para tarde (as
horas da \cref{prop:g5:healthy-choices:screens}, mais os custos da
\cref{prop:g8:protecting-the-nervous-system:sleep}). A defesa é a de
sempre: decidir as horas com antecedência e manter a última hora sem
tela.
\end{example}

\section{Prática: administrar a própria proteção}

\begin{method}[O manual do dono do sistema nervoso]\label{met:g8:protecting-the-nervous-system:manual}
A lista de cuidados do \cref{met:g5:healthy-choices:run}, reemitida
com os mecanismos:
\begin{enumerate}
  \item \emph{durma as horas} --- a consolidação e a faxina acontecem
        de noite; proteja a última hora das telas;
  \item \emph{faça orçamento do volume} --- o capital de receptores
        não volta a crescer: volume abaixo, pausas dentro, protetor
        auricular onde o som ruge;
  \item \emph{capacete} onde as quedas são rápidas
        (\cref{ex:g5:brain-in-command:helmet}) --- a fiação não se
        conserta como o \omterm{def:g3:bones-and-muscles:skeleton}{osso};
  \item \emph{nada de enganadores de junção} para um \omterm{def:g5:brain-in-command:system}{cérebro} em obras
        --- e o ``não'' decidido com calma e com antecedência (a frase
        de bolso do \cref{ex:g5:healthy-choices:answers}, agora com a
        justificativa completa dela);
  \item \emph{use a máquina} --- estudo, esporte, música: as junções
        se fortalecem com tráfego, e um circuito treinado é o
        patrimônio que todo outro hábito protege.
\end{enumerate}
\end{method}

\begin{example}[A ajuda existe e funciona]\label{ex:g8:protecting-the-nervous-system:help}
Para quem já estiver enrolado --- com o próprio hábito ou com o de um
amigo, com substâncias ou com telas: os endereços do
\cref{ex:g8:contraception:help} servem aqui também. Enfermeiros
escolares e médicos tratam a \omterm{prop:g5:healthy-choices:substances}{dependência} como aquilo que este capítulo
mostra que ela é --- uma condição de fiação, e não um veredito de
caráter --- e junções religadas podem se readaptar: mais difícil do
que nunca ter começado, mas longe de ser sem esperança. O pior plano é
o do pátio da escola: o silêncio.
\end{example}

\section{Exercícios}

\begin{exercise}[$\star$]\label{exo:g8:protecting-the-nervous-system:1}
O que o \omterm{prop:g1:healthy-habits:needs}{sono} faz nas junções, e quem mais precisa das horas dele?
\end{exercise}

\begin{exercise}[$\star$]\label{exo:g8:protecting-the-nervous-system:2}
Por que o som alto é uma questão de orçamento? O que é gasto, e qual é
o aviso de dose excessiva?
\end{exercise}

\begin{exercise}[$\star$]\label{exo:g8:protecting-the-nervous-system:3}
Onde todas as substâncias psicoativas agem, e de que três jeitos?
\end{exercise}

\begin{exercise}[$\star$]\label{exo:g8:protecting-the-nervous-system:4}
O que o álcool faz com o circuito --- diga três estações.
\end{exercise}

\begin{exercise}[$\star$]\label{exo:g8:protecting-the-nervous-system:5}
Enuncie o mecanismo da \omterm{prop:g5:healthy-choices:substances}{dependência} em duas frases.
\end{exercise}

\begin{exercise}[$\star$]\label{exo:g8:protecting-the-nervous-system:6}
Por que o ciclo da tela pertence a este capítulo, mesmo sem molécula
nenhuma?
\end{exercise}

\begin{exercise}[$\star\star$]\label{exo:g8:protecting-the-nervous-system:7}
Explique o paradoxo da revisão com a
\cref{prop:g8:protecting-the-nervous-system:sleep}.
\end{exercise}

\begin{exercise}[$\star\star$]\label{exo:g8:protecting-the-nervous-system:8}
Por que as mesmas doses custam mais a um \omterm{def:g5:brain-in-command:system}{cérebro} \omterm{def:g3:stages-of-life:stages}{adolescente} que a um
adulto? Dois mecanismos (\cref{rem:g8:puberty:moods},
\cref{rem:g8:protecting-the-nervous-system:addiction}).
\end{exercise}

\begin{exercise}[$\star\star$]\label{exo:g8:protecting-the-nervous-system:9}
Ligue os metros do \cref{ex:g5:brain-in-command:driver} ao
amortecimento do álcool: o que aumenta, e por que a lei sobre dirigir
alcoolizado é uma lei sobre \omterm{def:g5:brain-in-command:reaction}{tempo de reação}?
\end{exercise}

\begin{exercise}[$\star\star$]\label{exo:g8:protecting-the-nervous-system:10}
Justifique de novo três itens do manual do dono, com os mecanismos
deles no nível das junções.
\end{exercise}

\begin{exercise}[$\star\star$]\label{exo:g8:protecting-the-nervous-system:11}
``Uma condição de fiação, e não um veredito de caráter.'' O que essa
mudança de enquadramento muda sobre procurar ajuda --- e sobre oferecer
ajuda a um amigo?
\end{exercise}

\begin{exercise}[$\star\star\star$]\label{exo:g8:protecting-the-nervous-system:12}
Escreva o argumento de encerramento do ano: a partir das fendas
ajustáveis da \cref{def:g8:nervous-communication:synapse}, derive tanto
a maior força do \omterm{def:g5:brain-in-command:system}{cérebro} quanto as vulnerabilidades características
dele --- aprendizado e \omterm{prop:g5:healthy-choices:substances}{dependência} como um mecanismo só, lido duas
vezes.
\end{exercise}

\section{Problema: A noite de prevenção}

\begin{problem}[{Problema de fim de semana --- a noite de pais e
estudantes da escola, roteirizada com
mecanismos}]\label{pb:g8:protecting-the-nervous-system:1}
A turma prepara a noite de prevenção da escola: quatro estações ---
\omterm{prop:g1:healthy-habits:needs}{sono}, som, substâncias, telas --- cada uma com uma demonstração e um
cartão de mecanismo.

\medskip\noindent\textbf{Parte I --- Os cartões das
estações.}
\begin{enumerate}
  \item Escreva o cartão do \omterm{prop:g1:healthy-habits:needs}{sono}: a demonstração são as marcas de
        duas pessoas voluntárias no teste da régua
        (\cref{met:g5:brain-in-command:ruler}), uma descansada contra
        uma que dormiu pouco. Preveja as marcas e dê o mecanismo de
        junção.
  \item Escreva o cartão do som: um medidor de decibéis, um par de
        protetores auriculares e a frase ``você não consegue fazer
        elas crescerem de novo''. Complete o mecanismo.
  \item Escreva o cartão das substâncias: um diagrama de \omterm{def:g8:nervous-communication:synapse}{sinapse},
        três setas com os rótulos imitar, bloquear e forçar. Atribua
        o álcool, a nicotina e um veneno do
        \cref{exo:g8:nervous-communication:11} às setas deles.
  \item Escreva o cartão das telas: sem molécula nenhuma --- então o
        que a demonstração do celular com notificações compartilha
        com a estação das substâncias?
\end{enumerate}

\medskip\noindent\textbf{Parte II --- As perguntas da
plateia.}
\begin{enumerate}[resume]
  \item Um pai: ``a gente bebia na idade deles e ficou bem''.
        Responda com a assimetria do canteiro de obras --- o que é
        diferente aos catorze anos?
  \item Uma estudante: ``um cigarro não religa nada''. Responda com a
        \cref{rem:g8:protecting-the-nervous-system:addiction} e a
        \cref{rem:g5:healthy-choices:never} --- qual é o preço de
        verdade da primeira dose?
  \item Uma mãe: ``as telas não são só a nova televisão?''. Distinga
        os esquemas: o que faz dos prêmios pequenos e imprevisíveis o
        treinador mais duro?
  \item Um estudante, baixinho: ``e se alguém já não consegue
        parar?''. Dê a resposta do
        \cref{ex:g8:protecting-the-nervous-system:help}, com os
        endereços incluídos.
\end{enumerate}

\medskip\noindent\textbf{Parte III --- A fala de
encerramento.}
\begin{enumerate}[resume]
  \item A fala abre com a malabarista e fecha com a \omterm{prop:g5:healthy-choices:substances}{dependência} ---
        ``a mesma capacidade de ajuste, lida duas vezes''. Redija esse
        argumento em quatro frases.
  \item Depois ela entrega a cada estudante a frase de bolso do
        \cref{ex:g5:healthy-choices:answers}. Explique, com mecanismo
        incluído, por que a frase precisa ser composta \emph{agora}.
  \item O cartaz da noite mostra os cinco itens do manual do dono.
        Ponha esses itens em ordem para a semana de alguém de catorze
        anos, do mais protetor ao menos, e defenda a sua primeira
        escolha.
  \item Encerre a noite numa frase: o que está sendo protegido, e por
        que os catorze anos ficam no meio da década decisiva?
\end{enumerate}
\end{problem}
